\specialchapter{致读者}

\begin{enumerate}
\def\labelenumi{\arabic{enumi}.}
\item
  本丛书（其目录载于第 ix 页与第 x 页）从数学的起点讲起，并给出完整的证明。原则上，它不要求读者具备任何特定的数学知识，而只要求对数学推理有一定的熟悉程度以及一定的抽象思维能力。尽管如此，它尤其面向那些至少已很好掌握大学数学课程第一年或第二年内容的读者。
\item
  我们选定的阐述方法是公理化且抽象的，并且通常从一般过渡到特殊。这一选择是由本论著的主要目的决定的，该目的就是为整个现代数学主体提供坚实的基础。为此，熟悉数量相当可观的非常一般的概念和原理是必不可少的。此外，证明的要求使题材有一个严格固定的顺序。因此，除非读者已经具备相当广泛的数学知识，某些考虑的效用是不会立即显现出来的；在这种情形下，他必须有耐心把判断悬置起来，等待时机的到来。
\item
  为了减轻这一不利之处，我们经常在正文中插入一些例子，它们涉及读者可能已经知道、但在本丛书中尚未讨论过的事实。这类例子总是置于两个星号之间：*···*。大多数读者无疑会发现这些例子有助于理解正文，并且即使初读也不愿略去它们。当然，从纯逻辑的观点看，略去它们并无任何不便。
\item
  本丛书分为若干卷（这里称为“书”）。一般地，正文中的每一陈述都只把前面各卷中已讨论过的结果视为已知；前六书编有号码。这条规则在每一书内部都成立，但为阐述方便起见，这些书不再按连贯的顺序排列。在每书（或其各章）的开头，读者会找到关于该书与其他各书之间逻辑关系的精确说明，从而他能够确信不存在任何恶性循环。
\end{enumerate}

\pandocbounded{\includegraphics[keepaspectratio]{images/part001_8f5d44fd.jpg}}

\begin{enumerate}
\def\labelenumi{\arabic{enumi}.}
\setcounter{enumi}{4}
\tightlist
\item
  每一章的逻辑框架由该章的定义、公理和定理组成。这些是日后使用时主要需要记住的部分。较次要的结果以及那些容易由定理推出的结果被称为“命题”、“引理”、“推论”、“注记”等。初读时可以略去的那些内容以小号字体排印。对特别重要定理的评注偶尔以“附释”之名出现。
\end{enumerate}

为避免乏味的重复，有时引入仅在某一章或某一章的某一节内有效的记号或缩写是方便的（例如，在只讨论交换环的一章中，“环”一词总是指“交换环”）。这类约定总是被明确指出，一般出现在使用它的那一章的开头。

\begin{enumerate}
\def\labelenumi{\arabic{enumi}.}
\setcounter{enumi}{5}
\tightlist
\item
  正文中的某些段落旨在预先提醒读者以防严重的错误。这些段落在页边用符号
\end{enumerate}

Z（“危险弯道”）。

\begin{enumerate}
\def\labelenumi{\arabic{enumi}.}
\setcounter{enumi}{6}
\item
  习题的目的，既在于使读者能确信自己已消化正文，也在于使他注意到那些在正文中没有位置但仍有意义的结果。最难的习题带有符号 ¶。
\item
  一般地说，我们沿用通行的术语，除非看起来有充分理由偏离它。
\item
  我们特别注意始终使用严格正确的语言，同时不牺牲简单性。在正文中，我们尽可能指出语言的滥用；没有这类滥用，任何数学文本都有流于迂腐之虞，更不用说不可读。
\item
  由于正文原则上是对一个理论的教条式阐述，它一般不包含对文献的引用。书目式引用汇集于历史注记中，通常位于每章之末。这些注记在适当之处还包含关于该理论未解决问题的提示。
\end{enumerate}

跟随每个历史注记之后的参考书目，一般只包含那些在所论理论的发展过程中具有最重要意义的书籍和原始论文。它绝无自诩完备之意；特别是，仅用于确定优先权问题的引用几乎总是被略去。

至于习题，我们一般并不认为指出它们的来源是值得的，因为它们取自许多不同的来源（原始论文、教科书、习题集）。

11. 对本丛书某一部分的引用按如下方式给出：\par

\begin{enumerate}
\def\labelenumi{\alph{enumi})}
\item
  如果引用的是同一节中给出的定理、公理或定义，则按其编号引用。
\item
  如果它们出现在同一章的另一节中，则在引用中还给出该节。
\item
  如果它们出现在同一书的另一章中，则给出章和节。
\item
  如果它们出现在另一书中，则先按标题给出该书。
\end{enumerate}

结果摘要用字母 R 引用：例如《集合论》R 表示“《集合论》的结果摘要”。

